% Cumulative addition. Insert after the existing Sturmian suspension.
% Preserve the complete text and figure fig:holonomia-esturmiana.
\subsection{Nonadic connection, holonomy, and monodromic representation}
\label{mono-ampl:conexion}

The already constructed regional connection is now studied through three objects:
transport of the complete state, the representation of its turns, and the
suspension of its itineraries. The APP--TRIT--TPK operations and the
initial catalog retain their preceding definitions. The necessary
distinction here concerns what each representation retains of the
same history. Definition~\ref{mono-ampl:def-vuelta} preserves the admissible
extensions; Propositions~\ref{mono-ampl:vueltas}
and~\ref{mono-ampl:representacion} distinguish their iteration and
the bilateral representation of the index.

At prolongation index \(k\), a state chart preserves
\[
 \widehat x_k=(B_k,C_k,p_k,c_k,\Sigma_k,W_k,g_k,m_k).
\]
Here \(B_k\) brings together the three regional blocks, \(C_k\) is the cylinder,
\(p_k\) the prefix, \(c_k\) the carry record, \(\Sigma_k\) the boundary
signature, and \(W_k\) the incidence information. The phase is
\(g_k=1+[(k-1)]_9\); \(m_k\) counts complete turns. The reduction
\(\beta_k=[m_k]_{12}\) preserves the dodecaphase position, whereas the
integer \(m_k\) distinguishes its successive turns.

\begin{definition}[Transport through one nonadic turn]
\label{mono-ampl:def-vuelta}
Let \(\mathcal R_k\subseteq\widehat X_k\times\widehat X_{k+1}\)
be the admissible-extension relations of the TPK. The realization of
holonomy at level \(k\) is the relational composition
\[
 \Gamma_{9,k}=\mathcal R_{k+8}\circ\cdots\circ\mathcal R_k.
\]
For \(s\geq1\), its prolongation to \(s\) turns is
\[
 \Gamma_{9s,k}=
 \Gamma_{9,k+9(s-1)}\circ\cdots\circ\Gamma_{9,k}.
\]
\end{definition}

The relational formulation preserves admissible continuations. On
an individualized history, each arrow leads to the state actually
prolonged; the isolated visible block may still admit more than one
continuation. This difference is precisely the function of memory and
boundary conditions in determining the trajectory.

\begin{proposition}[Iteration of return and accumulated memory]
\label{mono-ampl:vueltas}
In every admissible composition of \(s\) turns,
\[
 g_{k+9s}=g_k,\qquad m_{k+9s}=m_k+s,\qquad
 \beta_{k+9s}=\beta_k+s\pmod{12}.
\]
\end{proposition}
\begin{proof}
The one-turn rule preserves phase and increases the counter by one
unit. Applying it to the already transported state repeats exactly those
two operations. Induction gives the first two equalities; reducing the
second modulo \(12\) gives the third.
\end{proof}

The counter thus provides an integer witness to the difference between
turns. After \(12\) turns the pair of finite phases is restored, but
the counter increases by \(12\). The nonsplit extension
\eqref{eq:extension108} expresses the composition of the calendar through its
cocycle; the trajectory additionally retains the integer lift of that datum.
Conservation here means transport of previous contributions,
not immobility of every coordinate.

\subsubsection{Bilateral representation of the return index}

The shift \(T\) in \eqref{eq:weyl-relojes} admits an
expression in phase and turn coordinates. We number phases from \(0\) to
\(8\), by shifting the origin relative to \(g_k\), and
write \(k=9m+r\), with \(0\leq r<9\), on the integer line
covering the phase cycle. On
\(\mathcal H_{\rm ind}=\ell^2(\mathbb Z)\), the basis
\(|k\rangle=|r,m\rangle\) expresses that same unitary shift
\[
 T|r,m\rangle=
 \begin{cases}
 |r+1,m\rangle,&r<8,\\
 |0,m+1\rangle,&r=8.
 \end{cases}
\]
The inverse reduces \(r\) by one unit when \(r>0\), and sends
\(|0,m\rangle\) to \(|8,m-1\rangle\). Passage through the phase origin
therefore incorporates the exact quotient unit.

\Needspace{8\baselineskip}
\begin{proposition}[Representation of the monodromy of the cycle]
\label{mono-ampl:representacion}
Let \(S=T^9\). Then
\[
 S|r,m\rangle=|r,m+1\rangle,\qquad
 \mathbb Z\longrightarrow\mathcal U(\mathcal H_{\rm ind}),
 \quad s\longmapsto S^s
\]
is a faithful unitary representation of the turn group of the cycle.
\end{proposition}
\begin{proof}
Nine shifts preserve the remainder of division by \(9\) and
increase the quotient. The powers satisfy
\(S^{s+t}=S^sS^t\), and \(S^s|r,m\rangle=|r,m+s\rangle\) distinguishes
every integer \(s\). A permutation of an orthonormal basis is unitary.
\end{proof}

This representation corresponds to the covering of the cycle, whose turn
group is \(\mathbb Z\), and represents phase together with its counter. The
TPK state additionally contains the trajectory coordinates already described.
The bilateral extension admits negative indices as coordinates of the
covering; a trajectory starting from an initial condition uses
its admissible half-axis. This provides an operator inverse
of the index without turning every extension relation into a bijection of the
complete state.

With \(D_M|k\rangle=\lfloor k/9\rfloor|k\rangle\), the identity
\([D_M,S]=S\) expresses the memory increment on the dense domain of
finitely supported vectors. If \(\mathcal I|k\rangle=|-k\rangle\) and
\(P_0\) projects onto indices divisible by \(9\), one additionally obtains
\begin{equation}
 \mathcal I T\mathcal I=T^{-1},\qquad
 \mathcal I D_M\mathcal I=-D_M-(I-P_0).
 \label{mono-ampl:inversion-contador}
\end{equation}
Indeed, for \(k=9m+r\),
\(\lfloor-k/9\rfloor=-m\) if \(r=0\) and equals \(-m-1\) if \(r>0\).
Inversion thus preserves the boundary term of Euclidean division.

\subsection{Block clock and capacity clock}
\label{mono-ampl:relojes}

The index of an application of refinement and the number of available
decimal positions have related but different laws. To
express them, we reserve \(n\) for the ternary-block index and define
\[
 \rho=\log_{1000}729=\log_{10}9,\qquad
 \delta=1-\rho,\qquad
 \mathcal C_n=\lfloor(n+1)\rho\rfloor,\quad n\geq0.
\]
The letter \(\mathcal C_n\) here designates capacity, not the dodecaphase
record \(K\). At the canonical phase origin, the preceding mechanical
word provides
\[
 Z_n=\sum_{j=0}^{n}z_j(0)=\lfloor(n+1)\delta\rfloor.
\]

\begin{proposition}[Exact balance between capacity and vacancies]
\label{mono-ampl:balance-capacidad}
For \(n\geq0\) and, in the second identity, \(n\geq1\),
\[
 \mathcal C_n=n-Z_n,\qquad
 \mathcal C_n-\mathcal C_{n-1}=1-z_n(0).
\]
\end{proposition}
\begin{proof}
The irrationality of \(\delta\) implies
\(\lceil(n+1)\delta\rceil=\lfloor(n+1)\delta\rfloor+1\).
Then
\[
 \lfloor(n+1)(1-\delta)\rfloor
 =n+1-\lceil(n+1)\delta\rceil=n-Z_n.
\]
Subtracting two consecutive identities completes the proof.
\end{proof}

An event \(z_n=1\) adds a block to the history while maintaining capacity
\(\mathcal C_n\); an event \(z_n=0\) increases both counters. In
any interval, capacity increments and vacancies sum to
exactly its length. Reducing modulo \(q\), for \(q=9\) or
\(q=108\), gives
\[
 [\mathcal C_n]_q=[n]_q-[Z_n]_q.
\]
The uniform block phase and the capacity phase are related through
the integrated memory of vacancies. In particular, a retention of
capacity and a complete turn of the connection have different
mechanisms, although both may preserve a phase observation.

Monotone inversion of the clock gives a first-arrival time:
\[
 T_{\rm cap}(a)=\min\{n:\mathcal C_n=a\}
 =\left\lceil\frac a\rho\right\rceil-1,
 \qquad a\geq1.
\]
With \(\sigma=\rho^{-1}-1\), the time taken to gain \(b\) units
of capacity is
\[
 T_{\rm cap}(a+b)-T_{\rm cap}(a)
 =b+\lceil(a+b)\sigma\rceil-\lceil a\sigma\rceil.
\]
The difference of ceilings takes one of the two integers adjacent to
\(b\sigma\). Thus \(9\) or \(10\) blocks are obtained to gain
\(9\) units of capacity, and \(113\) or \(114\) to gain \(108\).
These are times of the inverse clock, measured from a first arrival; the
turn counter is always interpreted in the index to which it belongs.

\subsection{Symbolic suspension and aperiodic time crystal}
\label{mono-ampl:cristal}

Rotation coding combines phase periodicity, local recurrence,
and global aperiodicity. Let \(\mathcal X_\delta\subseteq\{0,1\}^{\mathbb Z}\)
be the closure, in the product topology, of the translates of the
bilateral word \(z_n(\theta)\). Denote by \(\sigma_{\rm sh}\) the
shift of its sequences. Retaining a nonadic phase, the step is
\[
 F(r,z)=([r+1]_9,\sigma_{\rm sh}z),\qquad
 (r,z)\in C_9\times\mathcal X_\delta.
\]

\begin{definition}[Symbolic aperiodic time crystal]
\label{mono-ampl:def-cristal}
The unit-roof suspension of the preceding system is the space
\[
 \mathcal S_\delta=
 \bigl((C_9\times\mathcal X_\delta)\times\mathbb R\bigr)/\sim,
 \qquad (y,t+1)\sim(Fy,t),
\]
with flow \(\Phi_s[y,t]=[y,t+s]\). This model of temporal order,
together with its vacancy coding and representation of the turn counter,
is called a \emph{symbolic aperiodic time crystal}.
\end{definition}

The adjective temporal designates the action of the flow and its return section;
aperiodicity characterizes its itineraries. The frequency of ones in every
sequence is \(\delta\): the uniform discrepancy bound of less than one,
obtained by telescoping in any window, persists on taking the closure.
A periodic sequence would have rational frequency. Therefore,
\(\mathcal X_\delta\) contains no periodic orbits. The suspension
also has no closed orbits: a return of the flow would require an integer
time and periodicity of the itinerary.

Finite words do recur. Itineraries of length \(m\)
arise from a partition of the circle by \(m+1\) points; each
interval is visited recurrently by the irrational rotation. Even
when times are restricted to a fixed phase, the step \(9\delta\) remains
irrational. This is why every word appears in all
\(9\) phases, and explains the already established complexities
\[
 p(m)=m+1,\qquad p_9(m)=9(m+1),\qquad p_3(m)=3(m+1).
\]
Recurrence of local configurations coexists with the absence of a
global period. Its topological entropy is zero because these complexities
grow linearly.

Realization through the observables \(V_9,W_\delta\) and shift
\(T\) expresses the same order through the covariances
\eqref{eq:weyl-relojes}. Frequencies of the uniform clock belong to
\(\frac19\mathbb Z+\delta\mathbb Z\), modulo \(1\). The capacity clock
additionally preserves the relation
\([\mathcal C_n]_9=[n]_9-[Z_n]_9\); it is not replaced by the uniform phase
when interpreting returns. The term \emph{time crystal} is thus
defined through a mathematical dynamical system. Its physical realization
requires specification of an energy dynamics, its temporal observables,
and its stability, data distinct from the present symbolic construction.

\subsection{Specular symmetry of the regions and double realization}
\label{mono-ampl:regiones}

Specular exchange acts on the regional block
\(B=(u^\pi,u^e,u^\varphi)\in(\mathbb F_3^6)^3\) by
\[
 \mathfrak cB=(u^e,u^\pi,u^\varphi).
\]
Its fixed axis and aggregate are, respectively, \(u^\varphi\) and
\(s=u^\pi+u^e\). This fixation refers to the involution within a fiber;
both coordinates may evolve during nonadic transport.
With \(q=u^\pi+u^e+u^\varphi\),
\(a=u^\pi+u^e-u^\varphi\), and \(d=u^\pi-u^e\), one recovers
\[
 u^\varphi=2(q-a),\qquad s=q-u^\varphi,\qquad
 u^\pi=2(s+d),\qquad u^e=2(s-d).
\]
The equalities hold componentwise in \(\mathbb F_3\), where
\(2^{-1}=2\). The mirror preserves \(q,a,s\) and reverses \(d\); this
antisymmetric coordinate individualizes the distribution that the aggregate preserves
without distinguishing it. The inversion \(\mathcal I\) of the temporal index and the
channel involution \(\mathfrak c\) consequently act on
different data.

Phase return has an explicit witness:
\[
 \begin{aligned}
 B_1&=(012222\mid121221\mid112202),\\
 B_{10}&=(101222\mid112221\mid112002).
 \end{aligned}
\]
Both positions have phase \(1\); the block, axis, and memory have
changed. Likewise, synchronization of the censuses at phases \(8\) and
\(9\) cancels their relative differences, while the diagonal component
passes from \((59,59,59)\) to \((43,43,43)\). Equality of a relative
observable does not exhaust the state determining it.

Double realization preserves that architecture. In Section
\ref{exc:doble-lectura}, the same prefix determines nested
cylinders on the one hand and a sequence of coded words with their frames on the other.
The first map passes to the limit by contraction of diameters; the second
satisfies \(r_{n,m}Q_n=Q_mp_{n,m}\) and defines \(Q_\infty\).
Figure \ref{mono-ampl:fig-regional} anticipates these two maps.
The subsequent incidence development constructs codes, lattices, and the
realization of Moonshine in their own domains; the automorphism group
acts on that realization, not on the digits through an implicit
identification.

% Figura acumulativa; destino figures/monodromia_regional_ampliacion.tex.
\begin{figure}[H]
\centering
\begin{tikzpicture}[x=1cm,y=1cm,>=Latex,font=\small,
  flow/.style={->,line width=.55pt},
  ann/.style={align=center,font=\footnotesize}]
\node[font=\bfseries] at (6.9,9.2)
  {Nonadic connection and regional coordinates};
\begin{scope}[shift={(3.4,6.8)}]
 \foreach \k in {1,...,9}{
  \pgfmathsetmacro{\ang}{90-40*(\k-1)}
  \coordinate (p\k) at (\ang:1.52);
  \fill (p\k) circle (1.4pt);
  \node[fill=white,inner sep=1.1pt,font=\footnotesize]
    at (\ang:1.83) {\(\k\)};
 }
 \foreach \k in {1,...,8}{
  \pgfmathtruncatemacro{\next}{\k+1}
  \draw[flow,shorten >=3pt,shorten <=3pt] (p\k)--(p\next);
 }
 \draw[flow,shorten >=3pt,shorten <=3pt] (p9)--(p1);
 \node[ann] at (0,.1) {\(g_k\in C_9\)\\\(\Gamma_{9,k}\)};
 \node[ann] at (0,-2.23) {One turn: \(g\mapsto g\), \(m\mapsto m+1\)};
\end{scope}
\node[ann,text width=5.25cm] at (10.1,7.75)
 {\(B_k=(u_k^\pi,u_k^e,u_k^\varphi)\)\\
 regional block in a fiber};
\node at (8.55,6.5) {\(u^\pi\)};
\node at (11.65,6.5) {\(u^e\)};
\draw[<->,line width=.65pt] (8.95,6.5)--node[above]{\(\mathfrak c\)}(11.25,6.5);
\draw[densely dashed,gray] (10.1,5.72)--(10.1,7.23);
\node[fill=white,inner sep=2pt] at (10.1,5.6){\(u^\varphi\)};
\node[ann] at (10.1,4.6)
 {Axis \(u^\varphi\) and aggregate \(u^\pi+u^e\)\\
 invariant under \(\mathfrak c\)};
\draw[flow] (5.45,6.8)--(7.25,6.8);
\draw[gray!60,line width=.3pt] (.55,4.02)--(13.25,4.02);
\node[ann,text width=12cm] (hist) at (6.9,3.47)
 {Compatible history: prefixes, cylinders, orientation, quotients,
 incidence frames, and memory};
\coordinate (fork) at (6.9,2.85);
\draw[flow] (hist.south)--(fork);
\draw[flow] (fork)--(3.6,2.2);
\draw[flow] (fork)--(10.3,2.2);
\node[ann,text width=5.6cm] at (3.6,1.65)
 {Archimedean realization\\
 \(I_{n+1}\subseteq I_n,\quad |I_n|=729^{-n}\)\\
 limit of regional prefixes};
\node[ann,text width=5.6cm] at (10.3,1.65)
 {Incidence realization\\
 \(r_{n,m}Q_n=Q_mp_{n,m}\)\\
 encoded words and compatible frames};
\node[ann,text width=5.6cm] at (3.6,.42)
 {\(\pi,\varphi,e\); dodecaphase determination of \(\alpha\)};
\node[ann,text width=5.6cm] at (10.3,.42)
 {Witt--Golay--Leech--\(V^\natural\)\\
 automorphisms of the exceptional realization};
\draw[flow] (3.6,1.03)--(3.6,.78);
\draw[flow] (10.3,1.03)--(10.3,.78);
\end{tikzpicture}
\caption{\fontsize{9}{11}\selectfont Nine phases of a connection and two realizations of a compatible
history. The cycle represents phase; each turn increments the counter.
The regional diagram represents the involution \(\mathfrak c\), which
interchanges closure and propagation and fixes the self-scaling axis in a fiber.
This invariance does not require the axis to remain constant during transport.
The two lower maps use the same prefix: one determines its numerical limit,
while the other preserves words and incidence frames. The maps of the second
realization are specified in Section \ref{exc:doble-lectura}; the exceptional
construction is developed afterwards.}
\label{mono-ampl:fig-regional}
\end{figure}


\clearpage
\subsubsection{Rectilinear representation of phases and return}
\label{mono-recta:desarrollo}

The representation on the segment \([0,10]\) preserves the positional
origin of APP and orders its nine interior marks. In this diagram,
the integers \(1,\ldots,9\) are phase representatives. Horizontal
position indicates the order of the charts; upper labels
specify regional operations. Figure
\ref{mono-recta:figura} thus brings together the initial line, fixation of the
self-scaling axis, the specular distribution of the two channels, and
return with memory. The axis and aggregate relations and the regional
census are made explicit below; return is governed by
Proposition~\ref{mono-ampl:vueltas}.

% Recovery of fig_nonadic_route_mini.tex from the academic synthesis.
% Labels designate regional operations; the sum is performed in F_3^6.
\begin{figure}[H]
\centering
\begin{tikzpicture}[x=1.12cm,y=1cm,font=\small,>=Latex,
  guide/.style={line width=.45pt,black!55}]
 \draw[line width=.8pt] (0,0)--(10,0);
 \foreach \k in {0,...,10}{
  \draw[guide] (\k,-.09)--(\k,.09);
  \node[below=2.5mm] at (\k,0) {$\k$};
 }
 \foreach \k in {1,...,9}{\fill (\k,0) circle (1.7pt);}
 \node[anchor=west,font=\footnotesize] at (0,3.05)
   {Generating segment $[0,10]$ and nine phase representatives};
 \draw[guide] (5,.12)--(5,1.7);
 \node[above,align=center,font=\footnotesize] at (5,1.7)
   {axis fixation\\$u^\varphi$};
 \draw[decorate,decoration={brace,amplitude=4pt},guide]
   (6,.65)--(8,.65);
 \node[above=5pt,align=center,font=\footnotesize] at (7,.65)
   {regional aggregate\\$u^e+u^\pi$};
 \draw[guide] (9,.12)--(9,1.7);
 \node[above,align=center,font=\footnotesize] at (9,1.7)
   {specular involution\\$\pi\leftrightarrow e$};
 \draw[->,line width=.7pt] (9,-.7)
   .. controls (9,-1.55) and (1,-1.55) .. (1,-.7);
 \node[fill=white,inner sep=3pt,font=\footnotesize] at (5,-1.23)
   {return $9\to1$;\quad $m\mapsto m+1$};
 \node[font=\footnotesize] at (5,-1.95)
   {$w_6\to w_{12}\to w_{18}\to w_{24}\to w_{30}\to R_{36}\to G_9$};
\end{tikzpicture}
\caption{Rectilinear representation of the nonadic connection. The mark
\(5\) locates the first strong saturation and the regional
self-scaling axis; the brace \(6\)--\(8\) identifies the ternary aggregate
fixed within each fiber; phase \(9\) selects the closure
orientation. Return preserves phase and increases memory. The
abscissae number operational positions, while the labels
\(u^\pi,u^e,u^\varphi\) designate blocks of \(\mathbb F_3^6\).}
\label{mono-recta:figura}
\end{figure}


Fixation of the axis is formulated in the fiber of a given boundary. If
\(q=u^\pi+u^e+u^\varphi\) and
\(a=u^\pi+u^e-u^\varphi\), then
\[
 u^\varphi=2(q-a),\qquad s=u^\pi+u^e=q-u^\varphi
 \quad\text{in }\mathbb F_3^6.
\]
The data \((q,a)\) determine the axis and aggregate; the ordered distribution
of \(s\) retains the specular freedom resolved by prolongation.
The fifth phase has a local solution and an extension, with boundary
datum \(c=111111\) and residual signature \(\rho_W=000\). This is the
meaning of its first strong saturation. The symbol \(\varphi\)
placed above the mark \(5\) identifies the regional axis entering
that operation.

The intermediate phases preserve this decomposition within each fiber,
while their boundary data evolve. The complete census of local
multiplicities and extensions of the distinguished branch is
\[
\begin{array}{c|rrrrrrrrr}
\text{phase}&1&2&3&4&5&6&7&8&9\\\hline
N_{\rm loc}&3&2&5&4&1&1&3&3&2\\
N_{\rm ext}&2&5&4&1&1&3&3&2&1
\end{array}
\]
Phase \(6\) combines local uniqueness and three prolongations; phase
\(7\) resolves a triple specular fiber; phase \(8\) synchronizes
the three regional censuses. Phase \(9\) preserves two local
distributions with the same axis and aggregate, one of which
prolongs. These counts express different functions within a
single connection and justify the marks on the line.

In particular, for phases \(6,7,8\), the regional sums are
respectively \(012100\), \(101001\), and \(112000\). The conservation
represented by the brace is an invariance under the specular
distribution within each fiber. Subsequent real evaluation applies the
readers of compatible cylinders to complete histories; the
internal sum represented in the figure belongs to the ternary domain.

Finally, the arc \(9\to1\) represents passage from one turn to the
next. The formulas of Section
\ref{mono-ampl:conexion} preserve the phase remainder and increase the
turn counter. The rectilinear diagram and the cyclic diagram in
Figure \ref{mono-ampl:fig-regional} therefore express two complementary
aspects of the same holonomy: the order of positions and
composition of the return on the state with memory.


\subsection{Temporal transport and tritic exponential composition}
\label{mono-ampl:euler}

The Euler relations of Section \ref{euler-trit-articulacion}
describe local realizations of the three regimes. For a fixed
signature \(\tau\), the generator satisfies \(J_\tau^2=-\tau I\);
algebraic conjugation sends \(J_\tau\) to \(-J_\tau\). Consequently,
\[
 \exp(sJ_\tau)\exp(tJ_\tau)=\exp((s+t)J_\tau),\qquad
 \overline{\exp(tJ_\tau)}=\exp(-tJ_\tau),
\]
and the product of an evolution with its conjugate is unity. In the elliptic
type this is realized through \(\cos^2t+\sin^2t=1\); in the hyperbolic type,
through \(\cosh^2t-\sinh^2t=1\); in the parabolic type,
\((I+tN)(I-tN)=I\), since \(N^2=0\). The same conservation law
admits three quadratic expressions.

The elliptic realization closes a turn; the nilpotent one preserves a
linear displacement; the hyperbolic one separates expansion and contraction with
unit determinant. At the already generated regional parameters,
\[
 \exp(\pi J_{\mathrm e})=-I,\qquad
 \exp(J_{\mathrm p})=I+J_{\mathrm p},\qquad
 \exp(2\log\varphi\,J_{\mathrm h})=F_{\rm av}^{\,2}.
\]
The evaluation \(\exp(I)=eI\) recognizes unity in the common
exponential law. The number \(e\) is not assigned exclusively to the parabolic regime.

Transport of a local realization by an isomorphism \(U\) preserves
these relations: if \(J'=UJU^{-1}\), the exponential series gives
\(U\exp(tJ)=\exp(tJ')U\). The identity transports the regime and its
norm. A transition between different quadratic types belongs, by contrast,
to the TPK selection of sheets and regimes; it is not a conjugation
changing the value of \(J^2\). The temporal product therefore preserves the
order of transitions and the information at their endpoints.

This articulation distinguishes closure of a local representation from
holonomy of a history. The identity \(\exp(2\pi J_{\mathrm e})=I\) may hold
at the same time as \(S|r,m\rangle=|r,m+1\rangle\). The first equality
restores the orientation of the regime; the second records another turn. The
discrete structure preserves both: local return and the accumulated
provenance of the state realizing it.
